\section{Retrievable Gradients and Dynamic Weight Editing}
\label{sec:regrad}

The parameter-efficient methods of Section~\ref{sec:peft} still modify weights permanently. A recent semi-parametric alternative seeks the quality of parametric updates without their cumulative cost, and serves here as the principal point of contrast for the retrieval-centric design adopted in this work.

Su et al. \parencite{su2026retrievablegradientscontinualposttraining} propose \textbf{Retrievable Gradients (REGRAD)} to address \emph{cumulative weight drift} in sequential continual post-training. Rather than updating the backbone, REGRAD pre-computes the LoRA gradient of each new document offline and stores it in a \emph{gradient bank}; at inference time, the relevant gradients are retrieved, temporarily added to the frozen weights, and then rolled back. This achieves parameter-level update quality while avoiding permanent drift. As Table~\ref{tab:regrad} shows, REGRAD attains the highest law-domain accuracy ($84.87\%$ with in-context learning) and the best overall average, all with zero weight drift, whereas standard continual pre-training degrades general-domain performance to $25.67\%$ due to severe drift.

\begin{table}[H]
\centering
\small
\caption{Average accuracy across domains for continual post-training methods on LLaMA-8B \parencite{su2026retrievablegradientscontinualposttraining}.}
\label{tab:regrad}
\begin{tabularx}{\textwidth}{X c c c c c}
\hline
\textbf{Method} & \textbf{General} & \textbf{Medical} & \textbf{Law} & \textbf{Overall} & \textbf{Drift?} \\
\hline
Direct baseline & 33.42 & 62.95 & 50.15 & 50.31 & No \\
Standard CPT & 25.67 & 62.03 & 48.20 & 48.18 & Yes \\
RAG (in-context) & 26.98 & 59.49 & 53.23 & 47.98 & No \\
Fine-tuned RAG & 37.17 & 66.84 & 70.33 & 61.34 & Yes \\
REGRAD & 38.43 & 65.61 & 80.50 & 64.36 & No \\
REGRAD + in-context & 45.06 & 67.74 & 84.87 & 67.26 & No \\
\hline
\end{tabularx}
\end{table}

Despite these results, REGRAD incurs high inference latency from applying gradients dynamically at runtime and requires maintaining a gradient-bank infrastructure. Crucially, like retrieval it avoids weight drift, but unlike retrieval it neither supports instantaneous knowledge updates nor offers a natural mechanism for intentional forgetting. These limitations---latency, infrastructure complexity, and difficult unlearning---are precisely why this work adopts retrieval-augmented generation as its primary mechanism rather than a gradient bank, while retaining REGRAD as the parametric-temporary point of comparison.
